\chapter{Поле реальних котушок}\label{ch:g09}

Стенд розд.~\ref{ch:g08} вносить збурення аналітично, з радіальною обвідною $g=\psiN^{m/2}$.
Вона росте до краю, як і належить вакуумному полю зовнішніх котушок. Але чи схожа вона на поле
\emph{справжнього} масиву RMP-котушок (resonant magnetic perturbation)? Реєстр візуалізації
записав це як здогадку VC1 з критерієм спростування «порівняння з полем I-котушок DIII-D за
Біо--Саваром» і вартістю «$\sim$день» \cite{RegViz}. Розділ розповідає, як здогадку спростовано
(VR14), що довелося побудувати натомість і яку помилку представлення поля ми зробили дорогою.
Та помилка~--- родичка R8 батьківського реєстру (розд.~\ref{ch:g07}).

\section{Біо--Савар, перевірений проти аналітики}\label{sec:g09-bs}

Масив I-котушок моделюється прямокутними «рамками» на циліндрі $R=R_{\text{coil}}$: два ряди по
6 котушок, струм у котушці $j$ зважено як $\cos(n\varphi_j)$ (\texttt{rmp\_coils.icoil\_array}).
Поле рахує сума внесків прямих відрізків:

\code{viz/src/tokviz/rmp_coils.py}{76}{99}{(\texttt{biot\_savart})}

Рядки 83--86: кожна петля~--- ламана з вершин \texttt{v}; відрізок від $\vect a$ до $\vect b$.
Рядки 87--92: точки обробляються порціями (\texttt{chunk}), щоб масив «точки $\times$ відрізки»
вмістився в пам'ять. Рядки 93--96: замкнена формула поля скінченного прямого відрізка,
$\dd\vect l\times\vect r_1\,(|\vect r_1|+|\vect r_2|)/[|\vect r_1||\vect r_2|(|\vect r_1||\vect r_2|+\vect r_1\cdot\vect r_2)]$.
Точки на самому провіднику (знаменник $\to0$) відкидаються. Рядок 98: множник $\mu_0I/4\pi$.

\begin{established}{VE27}
Поле на осі кругової петлі збігається з $\mu_0Ia^2/[2(a^2+z^2)^{3/2}]$ до 4 знаків при
$z/a=0;\ 0{,}5;\ 1;\ 2$ (тест \nolinkurl{test_biot_savart_matches_analytic_loop}). Масив із 6
котушок дає домінантну $n=3$ і першу бічну смугу $n=9$ на рівні 12\,\%. Саме таке аліасування
$n=3\pm6k$ має давати дискретний масив.
\end{established}

Дві помилки дорогою, обидві записані \cite{g07Campaign}. Перша (VE28): поле масиву спершу
зондували на магнітній осі при $Z=0$, де внески верхнього й нижнього рядів взаємно знищуються за
симетрією, і отримали абсурдні $4\cdot10^{-23}$~Тл. На краю плазми ($R=2{,}05$, $Z=0{,}3$) те саме
поле дає rms $2{,}8\cdot10^{-7}$~Тл/А. Друга: струми спочатку задавали бінарним знаком
$\operatorname{sign}(\cos n\varphi)$, а такий меандр сам породжує зайві гармоніки. Тепер у коді
стоїть синусоїдальне ваговання з коментарем, чому саме так (\texttt{rmp\_coils.py}, рядки 63--67).

\section{VC1: чи схожа обвідна на поле котушок}\label{sec:g09-vc1}

\paragraph{Метод.} Беремо нормальну до магнітної поверхні компоненту поля
$b_n=\vect B\cdot\grad\psi/|\grad\psi|$ і розкладаємо її в ряд Фур'є за $\thetastar$ і $\varphi$.
Модуль резонансної гармоніки $|b_{mn}(\psiN)|$ порівнюємо з тим самим числом для аналітичного
стенда. Амплітуду нормують, бо перевіряється лише \emph{форма} профілю.

\code{viz/src/run_coil_compare.py}{66}{82}{(\texttt{resonant\_harmonic})}

Рядок 71: точки поверхні рівномірно за $\thetastar$ і нормаль (функція \texttt{surface\_grid},
рядки 44--63). Рядки 76--78: $b_n$ на сітці $\varphi\times\thetastar$ для будь-якого поля, заданого
функцією. Тому одна процедура міряє і котушки, і стенд. Рядки 80--81: проєкція на
$\exp[-\mathrm i(m\thetastar-n\varphi)]$ дає комплексну амплітуду гармоніки.

Профіль стискаємо до одного числа~--- показника $p$ у $|b_{mn}|\propto\psiN^{p}$ на зовнішній
частині $\psiN\ge0{,}45$. Якщо обвідна чесна, для котушок має вийти $p$ поблизу
показника самого стенда:

\code{viz/src/run_coil_compare.py}{160}{171}{(\texttt{fit\_exponent})}

Рядок 168: лише скінченні додатні точки зовнішньої області. Рядок 171: нахил прямої в
логарифмічних координатах. Оскільки опубліковані розміри I-котушок у джерелах різняться,
порівняння повторено для шести варіантів геометрії (рядки 110--117 того самого файлу): номінал
$R_{\text{coil}}=2{,}184$~м, $z_c=0{,}754$~м, висота $0{,}5$~м, радіус $\pm10\,\%$, висота центру
$\pm25\,\%$ і «високі» котушки 0,8~м.

\begin{table}[ht]
\centering\small
\caption{Показник $p$ і кореляція профілів котушок зі стендом. Дані:
\texttt{coil\_compare.json}, звірено для путівника}
\label{tab:g09-exp}
\begin{tabular}{@{}lcccc@{}}
\toprule
& \multicolumn{2}{c}{$m/n=2/1$} & \multicolumn{2}{c}{$m/n=3/1$}\\
варіант & $p$ & кор. & $p$ & кор.\\
\midrule
номінал & $-0{,}08$ & 0,61 & 2,44 & 0,94\\
$R-10\,\%$ & 0,78 & 0,63 & 1,55 & 0,37\\
$R+10\,\%$ & $-0{,}34$ & 0,25 & 0,94 & 0,91\\
$z_c-25\,\%$ & $-0{,}93$ & $-0{,}45$ & 0,07 & 0,71\\
$z_c+25\,\%$ & 1,01 & 0,95 & 2,96 & 0,77\\
високі & $-0{,}48$ & 0,14 & 0,86 & 0,86\\
\midrule
стенд, $\psiN^{m/2}$ (та сама метрика) & 0,73 & --- & 1,23 & ---\\
\bottomrule
\end{tabular}
\end{table}

\begin{figure}[ht]
  \centering
  \includegraphics[width=\textwidth]{g09_coil_compare.jpg}
  \caption{Нормований радіальний профіль резонансної нормальної компоненти: шість варіантів
  масиву котушок (кольорові) проти обвідної стенда $\psiN^{m/2}$ (чорна штрихова). Жоден варіант
  котушок не монотонний назовні. Власний розрахунок, код:
  \texttt{tokamak-3d-viz/src/run\_coil\_compare.py}.}
  \label{fig:g09-coil}
\end{figure}

\begin{refuted}{VC1 $\to$ VR14}
«Обвідна $\psiN^{m/2}$ достатньо близька до поля реальних RMP-котушок.» Показник степеня для
котушок гуляє від $-0{,}93$ до $+1{,}01$ для $m=2$ (обвідна в тій самій метриці дає 0,73) і від
0,07 до 2,96 для $m=3$ (обвідна 1,23). Кореляція профілів коливається від $-0{,}45$ до $+0{,}95$.
Розкид задають радіус і висота масиву~--- саме ті параметри, які в публікаціях подають
по-різному. Жоден варіант котушок не монотонний назовні, тоді як $\psiN^{m/2}$ монотонна за
побудовою. \emph{Наслідок:} обвідна придатна для \emph{візуалізації}, але не для кількісних
тверджень про реальні RMP-експерименти.
\end{refuted}

Зауважте формулювання «придатна для візуалізації». Спростування не викидає інструмент, а звужує
його застосовність. Це записано в реєстрі, щоб рендер стенда ніхто не прочитав як модель
конкретного розряду з котушками.

\paragraph{Побічна перевірка: паразитні гармоніки (VC2 $\to$ VE26).} Ця сама кампанія закрила й
сусідню здогадку. Збурення містило \emph{лише} 2/1 і 3/1. По 44 орбіти засіяно в смузі
$\pm0{,}045$ навколо семи кандидатних резонансів, трасування на 150 обертів. Керовані ланцюги
знайдено ($W=0{,}0756$ і $0{,}0651$ проти теорії $0{,}0736$ і $0{,}0619$). Усі некеровані
резонанси (5/2, 1/1, 3/2, 7/2, 4/1) дали рівно $0{,}0000$ при порозі $0{,}004$. Тест мусив бути
\emph{кількісним}: майже будь-яке $(m,n)$ досяжне як цілочисельне биття, наприклад
$(1,1)=2(2,1)-(3,1)$. Тому вирішує не те, \emph{який} резонанс з'явився, а його \emph{розмір}.
Залишкова паразитна $m=1$ на рівні 1,3\,\% дала б на $q=1$ ланцюг завширшки $\approx0{,}0077$,
тобто в 1,9 раза вищий за поріг. \emph{Поправка 2026-09-29:} спершу тут стояло «$\approx0{,}0120$,
утричі вище порогу», пораховане старою формулою з помилкою $\sqrt q$ (розд.~\ref{ch:g08}, VR15).
Запас зменшився, вердикт лишився: виміряно рівно нуль.

\section{Другий бекенд: тримати \texorpdfstring{$\vect A$}{A}, а не \texorpdfstring{$\Bvec$}{B}}\label{sec:g09-backend}

Після VR14 стенд отримав \emph{паралельний} бекенд збурення з полем реального масиву
\cite{g09CoilBackend}. Аналітичний бекенд лишився без змін, і всі попередні результати
відтворюються (регресія 13/13). Викликати Біо--Савара в правій частині ОДР надто дорого (сотні
відрізків на кожне зерно й кожен крок), тому поле обчислюють один раз на сітці $(R,Z)$ при
наборі $\varphi$, розкладають за тороїдальними гармоніками і сплайнують.

\paragraph{Перша спроба: тримати $\Bvec$.} Виміряна відносна $|\grad\cdot\Bvec|$ виявилась такою:
\begin{center}\small
\begin{tabular}{@{}lcc@{}}
\toprule
& ядро ($\psiN<0{,}5$) & край ($\psiN<0{,}95$)\\
\midrule
сирий Біо--Савар & $5\cdot10^{-7}$ & $1{,}5\cdot10^{-7}$\\
сітка 65, 5 гармонік & $4{,}4\cdot10^{-4}$ & $\mathbf{2{,}12}$\\
сітка 129, 8 гармонік & $4{,}4\cdot10^{-5}$ & $4{,}1\cdot10^{-3}$\\
\bottomrule
\end{tabular}
\end{center}
На краю дивергенція перевищувала \emph{сам масштаб поля}: провідники стоять усередині
розрахункової області, і біля них поле майже сингулярне. Найнебезпечнішим тут був \emph{вигляд}
помилки, а не її величина. Поле, що не зберігає потік, збиває $\psi$ уздовж лінії й руйнує
збіжність Біркгофа (розд.~\ref{sec:g08-classifier}), і класифікатор чесно доповідає «хаос». На
3~кА вийшло 19 хаотичних орбіт із 34 і жодної захопленої смуги. Це виглядало як переконлива
фізика дискретного масиву з бічними смугами.

\paragraph{Виправлення (U3).} Тримати \emph{векторний потенціал}, розкладений за гармоніками й
сплайнований, а $\Bvec$ брати як його аналітичний ротор:

\code{viz/src/tokviz/coilfield.py}{104}{127}{(\texttt{CoilPerturbation.\_A\_and\_derivs})}

Рядки 110--114: для кожної компоненти $A_R,A_\varphi,A_Z$ накопичуються значення та три
похідні. Рядки 115--121: для гармоніки $n$ беруться сплайни коефіцієнтів $\cos n\varphi$ і
$\sin n\varphi$. Рядки 123--124: похідні за $R$ і $Z$~--- аналітичні похідні тих самих сплайнів.
Рядок 125: похідна за $\varphi$~--- точно, з ряду Фур'є.

\code{viz/src/tokviz/coilfield.py}{130}{146}{(\texttt{delta\_B}: ротор у циліндричних координатах)}

Рядки 133--135 (докстринг) і 143--145: три компоненти $\grad\times\vect A$ у $(R,\varphi,Z)$.
Мішані похідні сплайна комутують точно, тому $\grad\cdot(\grad\times\vect A)=0$ \emph{структурно},
хоч якою грубою є сітка. Рядки 139--140: точки поза сіткою притискаються до її краю. Рядок 142:
\texttt{gain} дозволяє сканувати струм без перебудови (поле лінійне за струмом).

\code{viz/src/tokviz/coilfield.py}{161}{194}{(\texttt{divergence\_report}: вимір, а не припущення)}

Рядки 168--179: випадкові точки всередині плазми ($\psiN\le0{,}95$) і випадкові $\varphi$.
Рядки 180--189: $\grad\cdot\Bvec=\frac1R\pa_R(RB_R)+\frac1R\pa_\varphi B_\varphi+\pa_ZB_Z$
центральними різницями з кроком $10^{-4}$. Рядки 190--194: нормування на характерний масштаб
$|B|/R$. Докстринг прямо каже, навіщо функція: це ворота, які \emph{перевіряють} бездивергентність,
замість того щоб її припускати. Тест \nolinkurl{test_coil_backend_is_divergence_free} вимагає
відносного rms $<10^{-3}$.

\begin{established}{U3 (\texttt{COIL-BACKEND.md} §2; номер дав путівник)}
Відносна $|\grad\cdot\Bvec|$: $3{,}19\cdot10^{-2}$, якщо тримати $\Bvec$, і
$\mathbf{2{,}50\cdot10^{-6}}$, якщо тримати $\vect A$. Перевірки: $\grad\times\vect A$ проти
Біо--Савара~--- $2{,}6\cdot10^{-11}$; Біо--Савар проти кругової петлі~--- до 4 знаків;
реконструкція поля проти сирого Біо--Савара всередині плазми~--- медіана 0,1\,\%.
\end{established}

\paragraph{Аналогія з R8.} У батьківському реєстрі R8~--- спроба перенести Tokamap на суму мод
\emph{підстановкою} $V'(T)$. Симплектичність зламалась, $\max|\det J-1|=1{,}7$
(розд.~\ref{ch:g07}). Тут той самий клас провалу: збереження потоку втрачено через представлення
«заради зручності», а не через фізику. Різниця лише в симптомі: там помилка проявилась числом
$\det J$, тут маскувалась під хаос. Ліки теж однакові: будувати поле як \emph{ротор} (у
аналітичного бекенда це $\grad\times(\psi\grad\varphi)$, тут $\grad\times\vect A$), щоб закон
збереження виконувався за побудовою, а не з точністю інтерполяції. Для студента це найкорисніший
урок розділу: якщо чисельна схема може порушити закон збереження, помилка вам не повідомить про
себе, вона \emph{видасть себе за фізику}.

Фізично бекенди різняться не косметично. Аналітичне збурення є функцією потоку, тож
$\delta\Btor\equiv0$. Поле котушок має $\delta\Btor$, а воно входить у знаменник рівнянь силової
лінії $\dd R/\dd\varphi=R\BR/\Btor$. Тому \texttt{FieldLines} переписано так, щоб працювати через
$\Bvec$. Масив, що жене $n=1$, дає гребінку $n=1\pm6k$; виміряна потужність у $\vect A$:
$n=1{:}\ 1{,}000$, $n=5{:}\ 0{,}126$, $n=7{:}\ 0{,}041$, $n=11{:}\ 0{,}029$.

\section{Що не вийшло: забагато хаосу (U1, U2)}\label{sec:g09-open}

\begin{center}\small
\begin{tabular}{@{}lcccc@{}}
\toprule
струм (лише $n=1$) & $\delta B/\Bpol$ & регулярні & острівні & хаотичні\\
\midrule
$3\cdot10^2$~А & $1{,}8\cdot10^{-4}$ & 20 & 0 & 10\\
$1\cdot10^3$~А & $5{,}9\cdot10^{-4}$ & 18 & 0 & 12\\
$3\cdot10^3$~А & $1{,}8\cdot10^{-3}$ & 17 & 0 & 13\\
\bottomrule
\end{tabular}
\end{center}
Частка хаосу майже не росте з амплітудою: 33\,\% $\to$ 43\,\% при десятикратному зростанні струму.
Це підлога, а не відгук на збурення. Аналітичний бекенд при $\delta B/\Bpol\sim10^{-3}$ дає
одиниці відсотків хаосу і чіткі захоплені смуги.

\begin{refuted}{U1 (\texttt{COIL-BACKEND.md} §4)}
«Хаос породжує гребінка бічних смуг $n=5,7,11,\ldots$» Прямий тест: лише $n=1$~--- 21 хаотична
орбіта з 30; повна гребінка~--- 26 з 30. Смуги погіршують картину, але причиною не є.
\emph{Примітка путівника:} документ не вказує, за яких струму й засіву зроблено цей тест. Таблиця
вище для «лише $n=1$» дає 10--13 хаотичних із 30, тож числа двох прогонів безпосередньо не
порівнянні. Висновок усередині одного прогону (21 < 26) це не зачіпає. Неузгодженість зареєстровано
2026-09-29 як відкрите питання VQ5 реєстру візуалізації \cite{RegViz}.
\end{refuted}

\begin{conjecture}{U2 (відкрите)}
Дві правдоподібні причини лишаються нерозділеними.
\emph{Фізична:} поле котушок має широкий полоїдальний спектр і жене \emph{всі} резонанси $n=1$
($q=1,2,3,4,5$) одночасно. П'ять ланцюгів, що перекриваються, руйнують острівні ядра при
значно меншому $\delta B/B$, ніж два. Це узгоджується з тим, що в реальних RMP-експериментах
край стохастизується за помірних струмів.
\emph{Числова:} залишкова похибка реконструкції поля (медіана 0,1\,\%, максимум 9\,\% біля
провідників) широкосмугова і діє як шум на кожному кроці. Вона занижує збіжність Біркгофа, і
класифікатор читає це як хаос.
\emph{Що розділить:} прогін із суттєво дрібнішою сіткою й більшим числом гармонік. Якщо підлога
хаосу впаде, причина числова; якщо ні~--- фізична.
\end{conjecture}

Поки U2 не розділено, бекенд котушок придатний для геометрії силових ліній і порівняння
профілів, але \textbf{не} для статистики островів і хаосу. Для статистики лишається аналітичний
бекенд.

\begin{practicum}[котушки]
\begin{enumerate}
\item \verb|python src/run_coil_compare.py --out data/coil_compare| відтворює
      табл.~\ref{tab:g09-exp} і рис.~\ref{fig:g09-coil}. Додайте свій варіант геометрії до
      списку \texttt{variants} (рядки 110--117) і подивіться, як зсувається $p$.
\item Розділіть U2 (задача з розд.~\ref{ch:g10}). Створіть \texttt{CoilPerturbation} з
      \verb|grid_shape=(257, 257)| і довшим \verb|n_keep|, перевірте \verb|divergence_report()| і
      порахуйте частку хаосу тим самим \texttt{analysis.classify} для трьох струмів таблиці.
      Для порівняння \verb|python src/run_coil_bake.py --n-only 1| дає базовий прогін.
\item Контрольне питання: чому в \texttt{divergence\_report} похідну $\pa_R(RB_R)$ беруть від
      добутку $RB_R$, а не від $B_R$ (рядок 188)?
\end{enumerate}
\end{practicum}
